% Part: methods
% Chapter: induction
% Section: inductive-definitions

\documentclass[../../../include/open-logic-section]{subfiles}

\begin{document}

\olfileid[tr]{mth}{ind}{idf}

\olsection{Tümevarımsal Tanımlar}

Mantıkta nesne türlerini çok sık \emph{tümevarımsal olarak} tanımlarız;
yani tanımlanacak türden bir nesne sayılmanın kurallarını belirler ve bu
kurallarla o türden eski nesnelerden o türden yeni nesnelerin nasıl elde
edileceğini açıklarız. Örneğin bir dilin terimleri ve !!{formula}s gibi
özel simge dizisi türlerini çoğu kez tümevarımla tanımlarız. Basit bir
örnek olarak, $\mathrm{a}$, $\mathrm{b}$, $\mathrm{c}$, $\mathrm{d}$ harflerinden,
~$\circ$ simgesinden ve $[$ ile $]$ köşeli parantezlerinden oluşan
``$[[\mathrm{c} \circ \mathrm{d}][$'', ``$[\mathrm{a}[]\circ]$'',
``$\mathrm{a}$'' ya da ``$[[\mathrm{a} \circ \mathrm{b}]\circ
\mathrm{d}]$'' gibi dizileri ele alın. İlk iki dizide bir şeylerin
``yanlış'' olduğunu muhtemelen hissedersiniz: ilkinde köşeli parantezler
hiç ``dengelenmez'' ve ``$\circ$''un, kendileri anlamlı olan ifadeleri
``bağlaması'' gerektiğini düşünebilirsiniz. Üçüncü ve dördüncü dizi daha
iyi görünür: her ``$[$'' için onu kapatan bir ``$]$'' vardır (herhangi
biri varsa) ve her $\circ$ için iki yanında, bir çift köşeli parantezle
çevrelenmiş ``düzgün'' ifadeler bulabiliriz.

Neyin ``düzgün terim'' sayıldığını kesin olarak belirtmek istiyoruz.
Öncelikle her harf tek başına düzgündür. Tek başına bir harften ibaret
olmayan her şey, $s$ ile $t$'nin kendileri düzgün olmak üzere ``$[t
\circ s]$'' biçiminde olmalıdır. Tersine $t$ ile $s$ düzgünse aralarına
bir $\circ$ koyup onları bir çift köşeli parantezle çevreleyerek yeni
bir düzgün terim oluşturabiliriz. Bu işlemleri, düzgün terimler
kümesini \emph{tanımlamak} için kullanabiliriz. Bu bir \emph{tümevarımsal
  tanım}dır.

\begin{defn}[Düzgün terimler]
  \emph{Düzgün terimler} kümesi tümevarımsal olarak şöyle tanımlanır:
  \begin{enumerate}
  \item Her $\mathrm{a}$, $\mathrm{b}$, $\mathrm{c}$, $\mathrm{d}$
    harfi bir düzgün terimdir.
  \item $s_1$ ile $s_2$ düzgün terimlerse $[s_1 \circ s_2]$ de düzgün
    bir terimdir.
  \item Başka hiçbir şey düzgün terim değildir.
  \end{enumerate}
\end{defn}

Bu tanım, bir şeyin ancak ve ancak (1) ve~(2) koşullarına göre sonlu
sayıda adımda kurulabiliyorsa düzgün terim sayıldığını söyler. İlk
adımda yalnızca tek başına harflerden oluşan bütün düzgün terimleri
kurarız:
\[
\mathrm{a}, \mathrm{b}, \mathrm{c}, \mathrm{d}
\]
İkinci adımda (2)'yi kurduğumuz terimlere uygularız. İki harfin bütün
bileşimleri için şunları elde ederiz:
\[
[\mathrm{a} \circ \mathrm{a}], [\mathrm{a} \circ \mathrm{b}],
[\mathrm{b} \circ \mathrm{a}], \dots, [\mathrm{d} \circ \mathrm{d}]
\]
Üçüncü adımda (2)'yi şimdiye dek kurduğumuz herhangi iki düzgün terime
yeniden uygularız. $[\mathrm{a} \circ [\mathrm{a} \circ \mathrm{a}]]$
gibi yeni bir düzgün terim---burada $t$, 1.~adımdaki $\mathrm{a}$ ve
$s$, 2.~adımdaki $[\mathrm{a} \circ \mathrm{a}]$'dır---ve $[[\mathrm{b}
\circ \mathrm{c}] \circ [\mathrm{d} \circ \mathrm{b}]]$ terimini elde
ederiz; ikincisi 2.~adımdaki $[\mathrm{b} \circ \mathrm{c}]$ ile
$[\mathrm{d} \circ \mathrm{b}]$ terimlerinden kurulmuştur.
Bu böyle sürer. (3). madde, bu yolla kurulmamış herhangi bir şeyin
düzgün terimler kümesine sızmasını engeller.

``Düzgün'' gelen her simge dizisinin bu tanıma göre düzgün olduğunu
henüz kanıtlamadığımıza dikkat edin. Bununla birlikte kurabildiğimiz
her şeyin gerçekten ``düzgün'' geldiği açık olmalıdır: köşeli
parantezler dengelidir ve $\circ$, kendileri düzgün olan parçaları
bağlar.

Tümevarımsal tanımların temel özelliği şudur: bütün düzgün terimler
hakkında bir şey kanıtlamak istiyorsanız tanım, ele almanız gereken
durumları size söyler. Örneğin $t$'nin bir düzgün terim olduğu
söylenirse, tümevarımsal tanım $t$'nin nasıl görünebileceğini söyler:
$t$ bir harf olabilir ya da bazı düzgün $s_1$ ve~$s_2$ terimleri için
$[s_1 \circ s_2]$ olabilir. (3). madde gereği tek olasılıklar bunlardır.

Tümevarımsal olarak tanımlanmış bir kümenin bütünü hakkındaki iddiaları
kanıtlarken tümevarımın kuvvetli biçimi özellikle önem kazanır. Örneğin
uzunluğu~$n$ olan her düzgün terimdeki $[$ sayısının~$< n/2$ olduğunu
kanıtlamak istediğimizi düşünün. Bu, bütün~$n$'ler hakkında bir iddia
olarak görülebilir: her $n$ için, uzunluğu~$n$ olan herhangi bir düzgün
terimdeki $[$ sayısı $< n/2$'dir.

\begin{prop}
  Herhangi bir $n$ için, uzunluğu~$n$ olan bir düzgün terimdeki $[$
  sayısı $< n/2$'dir.
\end{prop}

\begin{proof}
Bu sonucu (kuvvetli) tümevarımla kanıtlamak için aşağıdaki koşullu
iddianın doğru olduğunu göstermeliyiz:
\begin{quote}
  Her $l < k$ için, uzunluğu~$l$ olan herhangi bir düzgün terimde $<
  l/2$ tane $[$ varsa, uzunluğu~$k$ olan herhangi bir düzgün terimde
  $< k/2$ tane $[$ vardır.
\end{quote}
Bu koşulluyu göstermek için önbileşeninin doğru olduğunu varsayalım;
yani herhangi bir $l<k$ için, uzunluğu~$l$ olan düzgün terimlerin $<
l/2$ tane $[$ içerdiğini varsayalım. Bu varsayıma tümevarım hipotezi
deriz. Aynı şeyin uzunluğu~$k$ olan düzgün terimler için doğru olduğunu
göstermek istiyoruz.

Öyleyse $t$, uzunluğu~$k$ olan bir düzgün terim olsun. Düzgün terimler
tümevarımsal olarak tanımlandığından iki durumumuz vardır: (1)~$t$ tek
başına bir harftir ya da (2)~$t$, bazı düzgün $s_1$ ve~$s_2$ terimleri
için $[s_1 \circ s_2]$'dir.
\begin{enumerate}
\item $t$ bir harftir. O zaman $k = 1$ ve $t$'deki $[$ sayısı~$0$'dır.
$0 < 1/2$ olduğundan iddia geçerlidir.
\item $t$, bazı düzgün $s_1$ ve~$s_2$ terimleri için $[s_1 \circ
  s_2]$'dir. $l_1$,~$s_1$'in; $l_2$ de~$s_2$'nin uzunluğu olsun.
  O zaman $t$'nin uzunluğu~$k$, $l_1+l_2+3$'tür ($s_1$ ve $s_2$'nin
  uzunluklarına üç $[$, $\circ$, $]$ simgesi eklenir).
  $l_1+l_2+3$ her zaman $l_1$'den büyük olduğundan $l_1 < k$.
  Benzer biçimde $l_2 < k$. Bu, tümevarım hipotezinin $s_1$ ve~$s_2$
  terimlerine uygulanabileceği anlamına gelir: $s_1$'deki $[$ sayısı
  $m_1$, $< l_1/2$'dir;~$s_2$'deki $[$ sayısı $m_2$ ise $< l_2/2$'dir.

  $t$'deki $[$ sayısı,~$s_1$'deki $[$ sayısı ile~$s_2$'deki $[$
  sayısına~$1$ eklenmiş hâlidir; yani $m_1 + m_2 + 1$'dir. $m_1 <
  l_1/2$ ve $m_2 < l_2/2$ olduğundan şunu elde ederiz:
  \[
  m_1 + m_2 + 1 < \frac{l_1}{2} + \frac{l_2}{2} + 1 = \frac{l_1+l_2+2}{2} < \frac{l_1+l_2+3}{2} = k/2.
  \]
\end{enumerate}
Her durumda, $t$'deki $[$ sayısının (tümevarım hipotezine dayanarak)
$< k/2$ olduğunu gösterdik. Kuvvetli tümevarımla önerme çıkar.
\end{proof}

\begin{prob}
  Süperdüzgün terimler kümesini şöyle tanımlayın:
  \begin{enumerate}
  \item Her $\mathrm{a}$, $\mathrm{b}$, $\mathrm{c}$, $\mathrm{d}$
    harfi bir süperdüzgün terimdir.
  \item $s$ bir süperdüzgün terimse $[s]$ de öyledir.
  \item $s_1$ ve $s_2$ süperdüzgün terimlerse $[s_1 \circ s_2]$ de
    bir süperdüzgün terimdir.
  \item Başka hiçbir şey süperdüzgün terim değildir.
  \end{enumerate}
  Uzunluğu~$n$ olan bir süperdüzgün~$t$ terimindeki $[$ sayısının
  $\le n/2 +1$ olduğunu gösterin.
\end{prob}

\end{document}
